\section{Summary of Contributions}

This thesis developed and evaluated a two-stage, explainable machine learning framework for predicting automotive fuel efficiency and estimating deterministic CO\textsubscript{2} emissions from vehicle specification data, directly adapting the methodological template of Atique et al.~\cite{atique2026explainable} to an independently assembled dataset and modeling pipeline. The principal contributions of this work are:

\begin{itemize}
    \item A reproducible, end-to-end ML pipeline encompassing data cleaning (duplicate removal, IQR-based outlier exclusion, and median imputation for missing values), domain-informed feature engineering (eight engineered features capturing engine technology evolution, drivetrain configuration, vehicle class, and brand-tier effects), and a broad benchmark of \textbf{seventeen supervised regression algorithms} spanning linear (linear, ridge, lasso), tree-based ensemble (Random Forest, Gradient Boosting, ExtraTrees, AdaBoost, XGBoost, LightGBM, CatBoost, Decision Tree), instance-based (KNN, SVR), and neural (MLP) model families, together with three advanced modeling techniques (stacking ensemble, Bayesian XGBoost, and a tuned feedforward neural network).

    \item A controlled ablation study empirically isolating the marginal contribution of the eight engineered interaction features to predictive performance, disaggregated by model family. The study demonstrated that feature engineering selectively benefited linear and regularized models (Lasso: \(\Delta R^2 = +0.0181\); Ridge: \(\Delta R^2 = +0.0066\)) and CatBoost (\(\Delta R^2 = +0.0106\)) under tuned settings, while tree-based ensembles with native interaction capacity (XGBoost, LightGBM, Gradient Boosting) showed marginal negative deltas, confirming that the utility of explicit feature construction depends on the inductive biases of the underlying algorithm.

    \item A tri-method explainability analysis --- SHAP, LIME, and permutation importance --- applied to the best-performing CatBoost model (\(R^2 = 0.9260\)), together with a formal Spearman rank-correlation analysis quantifying agreement between the three methods. The analysis found strong convergence between SHAP and permutation importance (\(\rho = 0.857\)), confirming cylinders, drive\_fwd, and displacement--year interaction as the dominant global predictors of fuel efficiency, while the weaker LIME alignment (\(\rho = 0.407\)) reflected its fundamentally local scope, providing complementary instance-level transparency rather than contradictory global rankings.

    \item A deterministic CO\textsubscript{2} emissions estimation module, translating predicted fuel efficiency (combination MPG) into fuel-type-specific tailpipe CO\textsubscript{2} estimates using emission factors of \(23.1\,\text{gCO}_2/\text{km}\) per L/100km for gasoline and \(26.4\,\text{gCO}_2/\text{km}\) per L/100km for diesel, explicitly excluding upstream and lifecycle emissions to maintain a well-defined boundary condition for tailpipe-only regulatory comparison.

    \item A comparative discussion situating the findings of this independent study within the existing literature, and in particular directly against the results reported by Atique et al.~\cite{atique2026explainable}, allowing an assessment of the generalizability of that study's conclusions beyond its original 550-record dataset. The replication of identical top-three features (cylinders, drive\_fwd, displacement--year interaction), near-identical Spearman correlations between SHAP and permutation importance (\(\rho = 0.857\) in both studies), and consistent algorithmic superiority of CatBoost across independently executed pipelines provides strong evidence that these mechanistic relationships are structurally stable rather than dataset-specific artifacts, substantially strengthening the external validity of the methodological template and the practical recommendations derived therefrom.
\end{itemize}


\section{Answers to Research Questions}

Based on the experimental results presented in Chapter~4 and the interpretive discussion in Chapter~5, the four research questions that motivated this thesis are answered as follows:

\paragraph{RQ1 (Best-performing algorithm):}
CatBoost emerged as the best-performing model across all evaluated configurations, achieving a test \(R^2\) of 0.9260, an RMSE of 1.1759, and an MAE of 0.8147 under tuned hyperparameter settings. This performance was closely followed by the stacking ensemble (\(R^2 = 0.9182\), RMSE \(= 1.2359\), MAE \(= 0.8414\)), which ranked second overall, and SVR with tuned hyperparameters (\(R^2 = 0.9188\), RMSE \(= 1.2315\), MAE \(= 0.9499\)), which was competitive but slightly inferior to CatBoost on the MAE metric. CatBoost's superiority is attributed to its ordered boosting mechanism, which systematically processes categorical features (e.g., vehicle class, make, drive type, transmission) without requiring extensive one-hot encoding, and its robust handling of the dataset's categorical-heavy structure. The model also demonstrated the most favorable balance between predictive accuracy and generalization stability, with a controlled train-test gap of 0.110 and the highest mean cross-validation \(R^2\) of \(0.8641 \pm 0.0324\), confirming that its strong test performance reflects genuine generalization rather than overfitting.

\paragraph{RQ2 (Contribution of feature engineering):}
The controlled ablation study revealed that feature engineering contributed selective but meaningful improvements, with the benefit depending critically on model family. Under tuned hyperparameter settings, feature engineering produced the largest gains for Lasso (\(\Delta R^2 = +0.0181\)), CatBoost (\(\Delta R^2 = +0.0106\)), and ExtraTrees (\(\Delta R^2 = +0.0096\)), followed by Ridge (\(\Delta R^2 = +0.0066\)) and Random Forest (\(\Delta R^2 = +0.0016\)). Linear and regularized models consistently benefited from the engineered features because explicit interaction terms (e.g., displacement--year interaction, cylinders--displacement interaction) capture domain-relevant nonlinearities that these models cannot otherwise represent. CatBoost's positive response is attributed to its ability to exploit explicit interaction terms as direct splitting candidates within its ordered boosting framework. Conversely, XGBoost (\(\Delta R^2 = -0.0075\)), LightGBM (\(\Delta R^2 = -0.0110\)), and Gradient Boosting (\(\Delta R^2 = -0.0025\)) exhibited marginal negative deltas under tuned settings, confirming that their native capacity to construct higher-order interactions through sequential tree growth rendered explicit features partially redundant or mildly disruptive. Overall, the engineered features delivered reliable gains for linear, regularized, and CatBoost models while remaining architecturally neutral or neutral for tree-based ensembles with inherent interaction capacity.

\paragraph{RQ3 (Most influential features and cross-method agreement):}
Across all three XAI methods, \textit{cylinders}, \textit{drive\_fwd}, and \textit{displacement\_year\_interaction} emerged as the most consistently influential predictors of fuel efficiency. SHAP analysis ranked cylinders first (mean absolute SHAP value dominating all other features), with high cylinder counts generating strongly negative contributions that reduced predicted MPG, while drive\_fwd (front-wheel drive) produced strong positive contributions. Permutation importance corroborated this ranking, with cylinder permutation degrading model \(R^2\) by \(0.297 \pm 0.021\), followed by drive\_fwd (\(\Delta R^2 = 0.092 \pm 0.015\)) and displacement--year interaction (\(\Delta R^2 = 0.088 \pm 0.017\)). LIME local explanations for individual test instances confirmed that low cylinder count was the strongest driver of elevated fuel efficiency, while the absence of front-wheel drive imposed the primary downward penalty. Cross-method consistency analysis revealed a strong Spearman correlation between SHAP and permutation importance (\(\rho = 0.857\), \(p < 0.0001\)), confirming robust global feature importance agreement, while SHAP versus LIME (\(\rho = 0.407\), \(p = 0.0009\)) and permutation importance versus LIME (\(\rho = 0.325\), \(p = 0.237\)) showed weaker alignment, reflecting LIME's fundamentally local scope rather than methodological inconsistency.

\paragraph{RQ4 (Practical implications):}
The XAI-confirmed feature relationships carry distinct practical implications for stakeholders. For automotive manufacturers, the dominance of cylinder count and displacement as negative predictors supports continued investment in engine downsizing, turbocharging, and front-wheel-drive platforms, with the displacement--year interaction statistically detecting the industry-wide trend toward smaller, more efficient powertrains. For environmental regulators and policymakers, the granular feature contributions quantified by SHAP and permutation importance could inform specification-based CO\textsubscript{2} emission benchmarks that reflect actual engineering relationships, while instance-level LIME explanations provide rule-based transparency for regulatory audit and vehicle certification. For consumers, the XAI transparency supports informed purchasing decisions by revealing how specific attribute choices (e.g., selecting a four-cylinder front-wheel-drive vehicle over a six-cylinder all-wheel-drive counterpart) translate into measurable fuel economy improvements, enabling more deliberate trade-offs between vehicle class, drivetrain, and engine configuration aligned with efficiency preferences and environmental values.



\section{Limitations Recap}

As detailed fully in Section~5.6, this thesis is limited by its reliance on static, specification-based data rather than real-world driving telemetry; by dataset-specific class imbalances, particularly across fuel type (546 gasoline, 2 diesel, 2 electric); by the regional/market scope of the underlying US EPA-style dataset; by collinearity among engine-size-related predictors (\(r = 0.94\) between displacement and cylinders); and by the deterministic, tailpipe-only scope of the CO\textsubscript{2} estimation module, which excludes upstream and lifecycle emissions. These limitations are consistent in kind with those explicitly acknowledged by Atique et al.~\cite{atique2026explainable} for their own dataset, and should be borne in mind when generalizing the findings of this thesis beyond the population represented in its dataset. The interpretability of the predictive models, while enhanced through three complementary XAI methods, remains bounded by the static nature of the input features and cannot account for real-world consumption variability arising from driver behavior, traffic conditions, road gradient, or ambient temperature. Similarly, the exclusion of electric vehicle lifecycle considerations and the near-absence of diesel and hybrid records constrain the scope of powertrain-related conclusions to the gasoline-dominated vehicle population that constitutes the dataset.


\section{Future Work}

Building on the limitations identified above, and following the future-work directions outlined by Atique et al.~\cite{atique2026explainable}, Section~5, three priority extensions are recommended for future research:

\begin{enumerate}
    \item \textbf{Validation on real-world driving telemetry.} Future work should validate the specification-based framework developed in this thesis against real-world telemetry collected via OBD-II ports or IoT-enabled sensors, in order to test its adaptability under dynamic driving conditions and to quantify the gap between specification-based and telemetry-based fuel-efficiency estimates, following the direction pursued by Canal et al.~\cite{canal2024machine} and Haghshenas et al.~\cite{haghshenas2025fuel}. Such validation would reveal whether the static feature hierarchy identified in this study (cylinders, drive\_fwd, displacement--year interaction) remains dominant under real-world conditions, or whether dynamic variables such as acceleration patterns, road gradient, and traffic density introduce entirely new primary predictors that overwhelm the specification-based signals.

    \item \textbf{Expansion to hybrid and electric vehicle populations.} Given the class imbalance documented in Section~3.1 (546 gasoline, 2 diesel, 2 electric) and analogous to that reported in Atique et al.~\cite{atique2026explainable}, future work should expand the dataset to include greater representation of hybrid and electric vehicles, in order to test whether the feature-importance hierarchy identified in this thesis is stable across powertrain types, or whether a materially different set of predictors (e.g., battery capacity, hybridization level, regenerative braking efficiency) emerges as dominant for electrified vehicles. This extension is particularly important as the automotive sector accelerates its transition toward electrification and policymakers require robust, evidence-based models that capture the distinct efficiency determinants of alternative powertrains.

    \item \textbf{Exploration of advanced tabular architectures and hybrid XAI pipelines.} Future work should explore modern tabular deep-learning architectures such as TabNet and FT-Transformer as alternatives to the tree-based ensembles and standard MLP explored in this thesis. These architectures, which employ attention mechanisms and feature-wise transformations, may offer advantages for capturing complex, non-linear interactions in structured vehicle data while potentially providing more direct mechanisms for feature attribution. Additionally, future work should investigate hybrid XAI pipelines that combine SHAP's global consistency with LIME's local fidelity in a single, interactive explainability interface suitable for real-time use by fleet operators, policymakers, and automotive engineers, as recommended by Atique et al.~\cite{atique2026explainable}.
\end{enumerate}

Beyond these three priority extensions, the following additional directions emerged from the experiments conducted in this thesis and warrant further investigation:

\begin{itemize}
    \item \textbf{Regional transferability and cross-market validation.} Given the US EPA-style origin of the dataset used in this thesis, future work should validate the framework on European (WLTP-cycle) and Asian (JC08 or region-specific) datasets to assess whether the feature importance hierarchy is stable across markets with different vehicle compositions, powertrain distributions, and regulatory standards. Such cross-market validation would significantly strengthen the generalizability claims of the methodological template.

    \item \textbf{Full lifecycle CO\textsubscript{2} estimation.} While this thesis limited its CO\textsubscript{2} estimation to deterministic tailpipe emissions, future work should extend the framework to incorporate full lifecycle emissions, including upstream fuel production, electricity generation for EVs, battery manufacturing, and vehicle production. This extension would align the predictive framework with comprehensive sustainability assessment goals and provide policymakers with a more complete carbon accounting tool.

    \item \textbf{Integration of temporal dynamics.} The dataset employed in this thesis spans model years 2014--2024, but the models treat time as a static feature. Future work should explicitly model temporal dynamics using time-series architectures, such as LSTM or Transformer-based models, to capture the evolution of vehicle efficiency over time and to improve forecasting of future fleet-level emission trajectories under different regulatory scenarios.
\end{itemize}



\section{Concluding Remarks}

This thesis has demonstrated that a carefully designed, two-stage machine learning framework --- combining a broad algorithmic benchmark, empirically validated feature engineering, and a triangulated, multi-method explainability analysis --- can deliver both strong predictive accuracy and transparent, mechanistically grounded insight into the determinants of automotive fuel efficiency. By independently re-applying and extending the methodological template of Atique et al.~\cite{atique2026explainable} to a new dataset, this work provides an additional data point in support of the broader claim that engine architecture --- captured through cylinder count, displacement, and their interaction with drivetrain configuration and model year --- constitutes a robust, largely model-agnostic determinant of vehicle fuel efficiency, and that this determinant can be surfaced reliably through the convergence of independent, complementary explainability techniques. The replication of identical top-three features (cylinders, drive\_fwd, displacement--year interaction) and near-identical cross-method Spearman correlations across independent pipelines substantially strengthens the external validity of both the reference study's conclusions and the present thesis's findings, confirming that these mechanistic relationships are structurally stable rather than dataset-specific artifacts. As the automotive sector continues its gradual transition toward electrification, frameworks of this kind --- interpretable, specification-based, and readily retrainable on new regional or powertrain-diverse data --- offer a practical, low-cost complement to expensive physical testing and dynamic telemetry-based approaches, supporting evidence-based decision-making for manufacturers, regulators, and consumers alike during this transitional period. The future work directions outlined above provide a clear roadmap for extending this framework to real-world telemetry, diverse powertrain populations, and advanced interpretability interfaces, ultimately contributing to a more data-driven and transparent approach to transportation sustainability.